% STATUS: written, 20 October 2026. Needs a figure reference in paragraph 4
% once Figure 1 exists.

\section{Introduction}

A distributed database splits its data into pieces and spreads those pieces
across machines in different places. The pieces are usually called shards, or
fragments, and the two decisions that matter are where the boundaries between
them fall and which machine holds each one. Both decisions are made once, when
the database is designed. The designer fixes the boundaries in a schema, and a
function of the key, such as a hash, decides the placement. That function maps a
record to a machine without ever considering where the requests for that record
come from.

That arrangement works whilst the traffic stays where the designer expected it.
In a service with users on several continents, served by a database spread
across regions in the manner of Spanner~\cite{corbett2013spanner} or
CockroachDB~\cite{taft2020cockroachdb}, it does not. The geographic center of
the workload moves across the day as one region wakes and another sleeps, and
abruptly when something local draws attention to one region, so a placement
chosen for the morning is wrong by the evening. Every request that has to cross
an ocean to reach its data pays the round-trip time of that ocean, and a
transaction touching shards held in several regions pays it repeatedly,
because keeping it atomic means coordinating every region it touches before it
can commit. Latency therefore rises steadily as the mismatch between placement
and traffic grows.

One response is to move the data when the traffic moves. A simple rule
does this: in each decision window (the fixed interval at which placement is reconsidered), find the worst-placed shard (the one that would gain most by moving)
and move it towards the region reading it. We implement that rule as a baseline
and call it the \emph{follow-traffic heuristic}, or simply the heuristic. This paper asks whether
anything is gained by replacing it with a policy learned from experience. How much
there is to gain depends on how fast the workload moves. When
drift is slow, the heuristic is 16.9\,ms above the lowest latency any
placement of the same shards could reach (a bound we find by exhaustive search and call the oracle), and there is little worth learning.
When drift is fast, the heuristic sits 75.4\,ms above that bound, 47\% of
its own latency, partly because it reacts one window behind a target that has
already moved again. The room for a learned policy, bounded above by the
oracle, is therefore largest where the problem is hardest.

We build a simulator of a six-region cluster whose workload center drifts on a
daily cycle and spikes locally, and on it we state the placement problem as a
Markov decision process with 261 actions, each of which names the placement it
wants to reach. Against static hash partitioning, which never moves anything, a
random control, the follow-traffic heuristic, a co-location planner modeled on
Lion and the oracle, we measure how the gap between the heuristic and the
oracle widens with workload volatility (how fast the workload's center moves,
and how often and how sharply it spikes). On a real six-node CockroachDB
cluster, with wide-area latency injected between the nodes, we verify that
placement changes transaction latency by far more than the measurement noise,
which must hold for anything learned in simulation to matter on a real
cluster. We then train a Deep Q-Network agent. At volatility 4, over four
training runs with different random seeds (learner seeds), it is 13.8\% faster
than static hashing and moves a fifth of the data the heuristic moves, but it
stays 17\,ms slower than the heuristic, which keeps the lower objective cost at
the migration weight used in training (the price the reward puts on moving
data). Run unchanged on the cluster on four workload seeds, beside
CockroachDB's own rebalancer, two trained agents again lose to the heuristic on
every seed, second place goes to one of them on three of the four, and the
rebalancer, left to itself, barely moves anything.
Fig.~\ref{fig:architecture} shows how the simulator, the agent and the
cluster fit together, and which three things have to agree between them for a
policy trained in one to mean anything in the other.
